\documentclass[11pt]{article}
\usepackage[margin=1in]{geometry}
\usepackage{amsmath,amssymb,amsthm}
\usepackage{booktabs}
\usepackage{enumitem}

\newtheorem{theorem}{Theorem}
\newtheorem{proposition}[theorem]{Proposition}
\newtheorem{corollary}[theorem]{Corollary}
\newtheorem{remark}[theorem]{Remark}

\title{d5\_245: a valid explicit baseline for \(G1\) and the reduced three-class splice target}
\author{}
\date{2026-03-22}

\begin{document}
\maketitle

\section{Purpose}
The remaining graph-side object \(G1\) was previously treated as a broad search
for a corrected selector package.
The goal of this note is to cut that search down to its actual core.
We do two things:
\begin{enumerate}[label=\textup{(\arabic*)},leftmargin=2.8em]
\item exhibit a completely explicit valid baseline package, proved directly for
all moduli \(m>2\), and
\item compare the surfaced mixed witness \texttt{mixed\_008} to that baseline at
the finite coarse-class level.
\end{enumerate}
The outcome is that the mixed witness differs from the valid baseline on only
three coarse classes, and exact coarse-class search shows that support-restricted
repairs on only three or four colors cannot work.
So the live content of \(G1\) is no longer a wide selector-family search: it is
best read as a \emph{three-class full-color splice} problem.

\section{An explicit valid baseline package}
For \(x=(x_0,\dots,x_4)\in (\mathbb Z/m\mathbb Z)^5\), write
\[
\sigma(x):=x_0+x_1+x_2+x_3+x_4 \pmod m.
\]
For each color \(c\in\{0,1,2,3,4\}\), define
\[
B_c(x)=x+e_{d_c(x)},
\]
where
\[
d_c(x)=
\begin{cases}
 c+1 \pmod 5, & \sigma(x)=0,\\
 c-1 \pmod 5, & \sigma(x)=1,\\
 c, & \sigma(x)\notin\{0,1\}.
\end{cases}
\]
Thus the package is:
\begin{itemize}[leftmargin=2.4em]
\item on exact sum layer \(0\): rotate the color by \(+1\),
\item on exact sum layer \(1\): rotate the color by \(-1\),
\item on all exact sum layers \(2,3,\dots,m-1\): use the color itself.
\end{itemize}

\begin{theorem}[explicit baseline package]
\label{thm:explicit-baseline}
For every modulus \(m>2\), the five maps \(B_0,\dots,B_4\) form a valid
colorwise package:
\begin{enumerate}[label=\textup{(\arabic*)},leftmargin=2.8em]
\item each \(B_c\) is a permutation of \((\mathbb Z/m\mathbb Z)^5\), and
\item at every vertex the chosen outgoing directions are exactly
      \(\{0,1,2,3,4\}\).
\end{enumerate}
\end{theorem}

\begin{proof}
For the outgoing statement, fix \(x\).
If \(\sigma(x)=0\), then the direction list is
\[
(c+1)_{c=0}^4=(1,2,3,4,0),
\]
which is a permutation of \(\{0,1,2,3,4\}\).
If \(\sigma(x)=1\), the direction list is
\[
(c-1)_{c=0}^4=(4,0,1,2,3),
\]
and otherwise it is simply
\[
(c)_{c=0}^4=(0,1,2,3,4).
\]
So outgoing exhaustiveness is immediate.

Fix a color \(c\).
Partition the torus into exact sum slices
\[
S_t:=\{x:\sigma(x)=t\},\qquad t\in\mathbb Z/m\mathbb Z.
\]
Because every step adds one basis vector, \(B_c\) always sends \(S_t\) into
\(S_{t+1}\).
Moreover, on each exact slice \(S_t\), the chosen direction \(d_c(x)\) is
constant: it depends only on whether \(t=0\), \(t=1\), or \(t\notin\{0,1\}\).
Hence the restriction
\[
B_c\vert_{S_t}\colon S_t\to S_{t+1}
\]
is just translation by a fixed basis vector, so it is bijective.
Since the exact sum slices partition the whole torus and each slice maps
bijectively to the next slice, \(B_c\) is a permutation.
\end{proof}

\section{Coarse classes for the surfaced mixed witness}
For the surfaced \texttt{mixed\_008} rule, the checked local feature rows collapse
into the following coarse classes:
\[
L0,\ L1,\ L2^{(0)},\ L2^{(1)},\ L3^{(0,p=0)},\ L3^{(0,p=1)},\ L3^{(1,p=0)},\ L3^{(1,p=1)},\ L4+.
\]
Here the superscript on \(L2\) is the old-bit value, while the superscript on
\(L3\) records the layer-3 old-bit \(s_3\in\{0,1\}\) and the simple predecessor
flag \(p\in\{0,1\}\) for the mixed witness.

The surfaced \texttt{mixed\_008} anchor table and the explicit baseline table are:
\begin{center}
\begin{tabular}{@{}lcc@{}}
\toprule
coarse class & mixed\_008 anchor & baseline anchor \\
\midrule
\(L0\) & 1 & 1 \\
\(L1\) & 4 & 4 \\
\(L2^{(0)}\) & 0 & 0 \\
\(L2^{(1)}\) & 2 & 0 \\
\(L3^{(0,p=0)}\) & 0 & 0 \\
\(L3^{(0,p=1)}\) & 3 & 0 \\
\(L3^{(1,p=0)}\) & 3 & 0 \\
\(L3^{(1,p=1)}\) & 0 & 0 \\
\(L4+\) & 0 & 0 \\
\bottomrule
\end{tabular}
\end{center}
So the mixed witness differs from the valid baseline on exactly three coarse
classes:
\[
L2^{(1)},\qquad L3^{(0,p=1)},\qquad L3^{(1,p=0)}.
\]

\begin{corollary}[three-class splice target]
\label{cor:three-class-splice}
At the coarse-class level, \(G1\) is not a general selector search anymore.
It is the problem of splicing the closed mixed branch into the explicit valid
baseline along exactly the three classes
\[
L2^{(1)},\ L3^{(0,p=1)},\ L3^{(1,p=0)}.
\]
\end{corollary}

\section{Exact coarse-class support search}
The surfaced mixed witness has a stable list of 60 bad outgoing direction
patterns on checked moduli \(m=5,7,9,11,13\).
A necessary support test on those 60 patterns shows that any repair support must
have size at least three, and the only support triples that even pass this
necessary test are the cyclic families
\[
\{0,1,2\},\ \{0,1,4\},\ \{0,3,4\},\ \{1,2,3\},\ \{2,3,4\}.
\]
However, exact CSP search on the full observed coarse-class carrier (141 class
tuples from the checked moduli) gives a stronger result.

\begin{proposition}[exact coarse-class search certificate]
\label{prop:coarse-search}
On the checked coarse-class carrier extracted from
\(m=5,7,9,11,13\):
\begin{enumerate}[label=\textup{(\arabic*)},leftmargin=2.8em]
\item no support of size three admits a coarse-class outgoing-exhaustive
      correction,
\item no support of size four admits a coarse-class outgoing-exhaustive
      correction,
\item the explicit baseline of Theorem~\ref{thm:explicit-baseline} is a valid
      full-support coarse-class solution.
\end{enumerate}
\end{proposition}

\begin{proof}[Proof sketch]
This is the exact checked content of the accompanying script
\texttt{d5\_244\_G1\_coarse\_class\_baseline\_search.py}.
The script enumerates the 141 coarse class tuples actually observed on the
checked moduli, then solves the outgoing-exhaustiveness CSP for every support
subset of sizes three and four.
All such supports fail.
It also checks directly, on \(m=3,5,7,9,11,13\), that the explicit baseline
package of Theorem~\ref{thm:explicit-baseline} is outgoing-exhaustive and
incoming-Latin.
\end{proof}

\section{Interpretation for the manuscript}
This reduction materially changes how the remaining graph-side work should be
read.
The live task is not to keep searching over broad mode-switch families.
Instead, one already has:
\begin{enumerate}[label=\textup{(\arabic*)},leftmargin=2.8em]
\item a direct explicit valid baseline package \(B\),
\item a surfaced closed mixed branch, and
\item the knowledge that coarse repairs supported on only three or four colors
      cannot work.
\end{enumerate}
So the natural next move is to combine this baseline with the finite-defect splice
principle of \texttt{d5\_236}: identify the actual defect image of the closed
mixed branch relative to the baseline, then certify the splice on those three
coarse classes.

\end{document}
